% Source: https://artofproblemsolving.com/community/c3h1503324_trapezoid_geometry_questions
% Stable identity: aops-topic-1503324-q01
\begin{asy}
usepackage("amsmath");
size(6cm);
defaultpen(fontsize(10pt));
pair S0=(0,0),R=(6,0),P=(1.5,3),Q=(4.5,3);pair H=(3,0);
pair G=extension(P,H,S0,Q),X=extension(P,R,S0,Q);
draw(P--Q--R--S0--cycle,linewidth(0.9));
draw(P--H,blue);draw(S0--Q,red);draw(P--R,darkgreen);
dot(G);dot(X);dot(H);
label("$P$",P,NW);label("$Q$",Q,NE);label("$R$",R,SE);label("$S$",S0,SW);label("$H$",H,S);
label("$G$",G,1.6*dir(162));label("$X$",X,1.3*S);
label("\mbox{}",(min(currentpicture,user=true)+max(currentpicture,user=true))/2,black); // keeps the text after the figure black when LaTeX typesets the labels inline
\end{asy}

\textbf{$G$ is the midpoint of $SQ$.} Since $H$ is the midpoint of $RS$, $SH=HR=\dfrac{RS}{2}=PQ$. So $PQRH$ has a pair of opposite sides $PQ$ and $HR$ that are parallel and equal, hence it is a parallelogram and $PH\parallel QR$. In triangle $SQR$, the line through the midpoint $H$ of $SR$ parallel to $RQ$ meets $SQ$ at its midpoint; so $SG=GQ=\dfrac{SQ}{2}$.

\textbf{Position of $X$.} Triangles $XPQ$ and $XRS$ are similar (alternate angles, $PQ\parallel RS$), with ratio $\dfrac{PQ}{RS}=\dfrac{1}{2}$. Hence
\[
\dfrac{PX}{XR}=\dfrac{1}{2}\qquad\text{and}\qquad\dfrac{QX}{XS}=\dfrac{1}{2},
\]
so $QX=\dfrac{SQ}{3}$ and $SX=\dfrac{2SQ}{3}$.

\textbf{Conclusion.} $X$ lies between $G$ and $Q$, and $XG=SX-SG=\dfrac{2SQ}{3}-\dfrac{SQ}{2}=\dfrac{SQ}{6}$. Therefore $\dfrac{XG}{QX}=\dfrac{\dfrac{SQ}{6}}{\dfrac{SQ}{3}}=\dfrac{1}{2}=\dfrac{PX}{XR}$, which rearranges to
\[
PX\cdot QX=XR\cdot XG .\qquad\blacksquare
\]
